\textbf{Phần 1. Trắc nghiệm lựa chọn}
\setcounterpageref{ex}{0}

\begin{ex} %[5D1V4]
Chuyển phân số $\dfrac{21}{60}$ thành phân số thập phân có mẫu số là $100$ ta được
\choice
{$\dfrac{30}{100}$}
{\True $\dfrac{35}{100}$}
{$\dfrac{40}{100}$}
{$\dfrac{42}{100}$}
\loigiai{
Trước hết ta rút gọn phân số $\dfrac{21}{60}$ bằng cách chia cả tử và mẫu cho $3$:
\[
\dfrac{21}{60} = \dfrac{21 : 3}{60 : 3} = \dfrac{7}{20}.
\]
Sau đó nhân cả tử và mẫu với $5$ để đưa về mẫu số $100$:
\[
\dfrac{7}{20} = \dfrac{7 \times 5}{20 \times 5} = \dfrac{35}{100}.
\]
\textbf{Đáp án của câu hỏi là B}
}
\end{ex}

\begin{ex} %[5D1V4]
Chuyển phân số $\dfrac{54}{125}$ thành phân số thập phân có mẫu số là $1\,000$ ta được
\choice
{$\dfrac{216}{1\,000}$}
{$\dfrac{540}{1\,000}$}
{\True $\dfrac{432}{1\,000}$}
{$\dfrac{423}{1\,000}$}
\loigiai{
Ta có $1\,000 : 125 = 8$. Nhân cả tử số và mẫu số với $8$:
\[
\dfrac{54}{125} = \dfrac{54 \times 8}{125 \times 8} = \dfrac{432}{1\,000}.
\]
\textbf{Đáp án của câu hỏi là C}
}
\end{ex}

\begin{ex} %[5D1V4]
Tìm chữ số $x$ thích hợp biết: $\dfrac{x}{25} = \dfrac{36}{100}$.
\choice
{$x = 4$}
{$x = 8$}
{\True $x = 9$}
{$x = 12$}
\loigiai{
Ta có: $\dfrac{36}{100} = \dfrac{36 : 4}{100 : 4} = \dfrac{9}{25}$.\\
Do đó $\dfrac{x}{25} = \dfrac{9}{25} \implies x = 9$.\\
\textbf{Đáp án của câu hỏi là C}
}
\end{ex}

\begin{ex} %[5D1V4]
Một người đi bộ trong $24$ phút được $2\text{ km}$. Hỏi trong $1$ giờ, người đó đi được bao nhiêu phần mười của $10\text{ km}$?
\choice
{\True $\dfrac{5}{10}$}
{$\dfrac{4}{10}$}
{$\dfrac{6}{10}$}
{$\dfrac{8}{10}$}
\loigiai{
Đổi $1$ giờ $= 60$ phút.\\
Trong $1$ phút người đó đi được: $2 : 24 = \dfrac{1}{12}\text{ (km)}$.\\
Trong $60$ phút người đó đi được: $\dfrac{1}{12} \times 60 = 5\text{ (km)}$.\\
Quãng đường $5\text{ km}$ so với $10\text{ km}$ là: $\dfrac{5}{10}$.\\
\textbf{Đáp án của câu hỏi là A}
}
\end{ex}

\begin{ex} %[5D1V4]
Viết đại lượng $350\text{ kg}$ dưới dạng phân số thập phân với đơn vị tấn ta được
\choice
{$\dfrac{35}{100}\text{ tấn}$}
{$\dfrac{350}{100}\text{ tấn}$}
{\True $\dfrac{350}{1\,000}\text{ tấn}$}
{$\dfrac{35}{10}\text{ tấn}$}
\loigiai{
Vì $1\text{ tấn} = 1\,000\text{ kg}$ nên $350\text{ kg} = \dfrac{350}{1\,000}\text{ tấn}$.\\
\textbf{Đáp án của câu hỏi là C}
}
\end{ex}

\begin{ex} %[5D1V4]
Phân số nào sau đây \textbf{không thể} viết được dưới dạng phân số thập phân?
\choice
{$\dfrac{3}{8}$}
{\True $\dfrac{5}{14}$}
{$\dfrac{7}{25}$}
{$\dfrac{9}{40}$}
\loigiai{
Một phân số tối giản muốn viết được thành phân số thập phân thì mẫu số của nó khi phân tích ra thừa số nguyên tố chỉ được chứa các thừa số $2$ và $5$.\\
Mẫu số $14 = 2 \times 7$ chứa thừa số nguyên tố $7$ khác $2$ và $5$, do đó phân số tối giản $\dfrac{5}{14}$ không thể viết được dưới dạng phân số thập phân.\\
\textbf{Đáp án của câu hỏi là B}
}
\end{ex}

\begin{ex} %[5D1V4]
Giá trị của biểu thức $\dfrac{7}{10} + \dfrac{15}{100} + \dfrac{45}{1\,000}$ viết dưới dạng phân số thập phân có mẫu số $1\,000$ là
\choice
{\True $\dfrac{895}{1\,000}$}
{$\dfrac{67}{1\,000}$}
{$\dfrac{875}{1\,000}$}
{$\dfrac{950}{1\,000}$}
\loigiai{
Quy đồng về mẫu số chung $1\,000$:
\[
\dfrac{7}{10} = \dfrac{700}{1\,000}; \quad \dfrac{15}{100} = \dfrac{150}{1\,000}.
\]
Khi đó:
\[
\dfrac{700}{1\,000} + \dfrac{150}{1\,000} + \dfrac{45}{1\,000} = \dfrac{700 + 150 + 45}{1\,000} = \dfrac{895}{1\,000}.
\]
\textbf{Đáp án của câu hỏi là A}
}
\end{ex}

\begin{ex} %[5D1V4]
Trong vườn thú có tất cả $80$ con thú, trong đó có $28$ con thú ăn cỏ. Phân số thập phân có mẫu số là $100$ chỉ số con thú ăn cỏ trong vườn thú là
\choice
{$\dfrac{28}{100}$}
{\True $\dfrac{35}{100}$}
{$\dfrac{32}{100}$}
{$\dfrac{40}{100}$}
\loigiai{
Phân số chỉ số con thú ăn cỏ là $\dfrac{28}{80}$.\\
Rút gọn bằng cách chia cả tử và mẫu cho $4$:
\[
\dfrac{28}{80} = \dfrac{28 : 4}{80 : 4} = \dfrac{7}{20}.
\]
Đưa về phân số thập phân có mẫu số $100$:
\[
\dfrac{7}{20} = \dfrac{7 \times 5}{20 \times 5} = \dfrac{35}{100}.
\]
\textbf{Đáp án của câu hỏi là B}
}
\end{ex}

\begin{ex} %[5D1V4]
Tìm phân số thập phân ở giữa hai phân số $\dfrac{2}{5}$ và $\dfrac{3}{5}$ có mẫu số là $100$:
\choice
{$\dfrac{40}{100}$}
{$\dfrac{60}{100}$}
{\True $\dfrac{50}{100}$}
{$\dfrac{70}{100}$}
\loigiai{
Đưa hai phân số về mẫu số $100$:
\[
\dfrac{2}{5} = \dfrac{40}{100}; \quad \dfrac{3}{5} = \dfrac{60}{100}.
\]
Phân số nằm giữa $\dfrac{40}{100}$ và $\dfrac{60}{100}$ là $\dfrac{50}{100}$.\\
\textbf{Đáp án của câu hỏi là C}
}
\end{ex}

\begin{ex} %[5D1V4]
Một cửa hàng bán được $45\text{ m}$ vải trong tổng số một cuộn vải dài $75\text{ m}$. Phân số thập phân biểu thị số mét vải cửa hàng đã bán là
\choice
{\True $\dfrac{60}{100}$}
{$\dfrac{45}{100}$}
{$\dfrac{50}{100}$}
{$\dfrac{75}{100}$}
\loigiai{
Phân số chỉ số mét vải đã bán là $\dfrac{45}{75}$.\\
Rút gọn bằng cách chia cả tử và mẫu cho $15$:
\[
\dfrac{45}{75} = \dfrac{3}{5}.
\]
Chuyển thành phân số thập phân có mẫu số là $100$:
\[
\dfrac{3}{5} = \dfrac{3 \times 20}{5 \times 20} = \dfrac{60}{100}.
\]
\textbf{Đáp án của câu hỏi là A}
}
\end{ex}

\begin{ex} %[5D1V4]
Cho phân số $\dfrac{a}{16}$. Tìm số tự nhiên $a$ bé nhất khác $0$ sao cho phân số đó có thể viết được dưới dạng phân số thập phân có mẫu số là $100$.
\choice
{$a = 2$}
{\True $a = 4$}
{$a = 6$}
{$a = 8$}
\loigiai{
Giả sử $\dfrac{a}{16} = \dfrac{k}{100}$ với $k$ là số tự nhiên.\\
Suy ra $a = \dfrac{16 \times k}{100} = \dfrac{4 \times k}{25}$.\\
Để $a$ là số tự nhiên nhỏ nhất khác $0$ thì $k$ phải là bội số nhỏ nhất của $25$ khác $0$, tức $k = 25$.\\
Khi đó $a = \dfrac{4 \times 25}{25} = 4$.\\
Thử lại: $\dfrac{4}{16} = \dfrac{1}{4} = \dfrac{25}{100}$ (thỏa mãn).\\
\textbf{Đáp án của câu hỏi là B}
}
\end{ex}

\begin{ex} %[5D1V4]
Một bể cá chứa $600\text{ lít}$ nước. Người ta tháo ra $252\text{ lít}$ nước để vệ sinh bể. Phân số thập phân biểu thị lượng nước còn lại trong bể so với sức chứa của bể là
\choice
{$\dfrac{42}{100}$}
{\True $\dfrac{58}{100}$}
{$\dfrac{38}{100}$}
{$\dfrac{62}{100}$}
\loigiai{
Lượng nước còn lại trong bể là:
\[
600 - 252 = 348\text{ (lít)}.
\]
Phân số chỉ lượng nước còn lại là:
\[
\dfrac{348}{600}.
\]
Chia cả tử và mẫu cho $6$ để đưa về mẫu số $100$:
\[
\dfrac{348 : 6}{600 : 6} = \dfrac{58}{100}.
\]
\textbf{Đáp án của câu hỏi là B}
}
\end{ex}

\textbf{Phần 2. Trắc nghiệm đúng sai}
\setcounterpageref{ex}{0}

\begin{ex} %[5D1V4]
Cho các phân số sau: $M = \dfrac{3}{8}$, $N = \dfrac{18}{200}$, $P = \dfrac{33}{300}$. Xét tính đúng sai của các khẳng định sau:
\choiceTF
{\True Phân số $M$ viết thành phân số thập phân có mẫu số $1\,000$ là $\dfrac{375}{1\,000}$.}
{\True Phân số $N$ rút gọn thành phân số thập phân có mẫu số $100$ là $\dfrac{9}{100}$.}
{Phân số $P$ viết thành phân số thập phân có mẫu số $10$ là $\dfrac{11}{10}$.}
{\True Tổng $N + P$ có giá trị bằng $\dfrac{20}{100}$.}
\loigiai{
\begin{itemize}
    \item $M = \dfrac{3 \times 125}{8 \times 125} = \dfrac{375}{1\,000}$. Khẳng định a \textbf{Đúng}.
    \item $N = \dfrac{18 : 2}{200 : 2} = \dfrac{9}{100}$. Khẳng định b \textbf{Đúng}.
    \item $P = \dfrac{33 : 3}{300 : 3} = \dfrac{11}{100}$ (mẫu số là $100$, không phải $10$). Khẳng định c \textbf{Sai}.
    \item $N + P = \dfrac{9}{100} + \dfrac{11}{100} = \dfrac{20}{100}$. Khẳng định d \textbf{Đúng}.
\end{itemize}
\begin{itemize}
    \item a) \textbf{Đúng}.
    \item b) \textbf{Đúng}.
    \item c) \textbf{Sai}: $P = \dfrac{11}{100}$.
    \item d) \textbf{Đúng}.
\end{itemize}
}
\end{ex}

\begin{ex} %[5D1V4]
Một thư viện trường học có $1\,000$ cuốn sách gồm: $450$ cuốn sách giáo khoa, $320$ cuốn truyện tranh, còn lại là sách tham khảo. Xét tính đúng sai của các khẳng định sau:
\choiceTF
{\True Phân số thập phân biểu thị số sách giáo khoa là $\dfrac{45}{100}$.}
{Phân số thập phân biểu thị số truyện tranh là $\dfrac{16}{100}$.}
{\True Số sách tham khảo trong thư viện là $230$ cuốn.}
{\True Phân số thập phân biểu thị số sách tham khảo là $\dfrac{230}{1\,000}$.}
\loigiai{
\begin{itemize}
    \item Sách giáo khoa: $\dfrac{450}{1\,000} = \dfrac{45}{100}$. Khẳng định a \textbf{Đúng}.
    \item Truyện tranh: $\dfrac{320}{1\,000} = \dfrac{32}{100} 
e \dfrac{16}{100}$. Khẳng định b \textbf{Sai}.
    \item Sách tham khảo: $1\,000 - (450 + 320) = 230$ (cuốn). Khẳng định c \textbf{Đúng}.
    \item Phân số chỉ sách tham khảo: $\dfrac{230}{1\,000}$. Khẳng định d \textbf{Đúng}.
\end{itemize}
\begin{itemize}
    \item a) \textbf{Đúng}.
    \item b) \textbf{Sai}: Phân số là $\dfrac{32}{100}$.
    \item c) \textbf{Đúng}.
    \item d) \textbf{Đúng}.
\end{itemize}
}
\end{ex}

\textbf{Phần 3. Trả lời ngắn}
\setcounterpageref{ex}{0}

\begin{ex} %[5D1V4]
Chuyển phân số $\dfrac{63}{150}$ thành phân số thập phân có mẫu số là $100$. Tử số của phân số thập phân đó là bao nhiêu?
\shortans{42}
\loigiai{
\textbf{Đáp số:} $42$\\
Chia cả tử số và mẫu số cho $3$:
\[
\dfrac{63}{150} = \dfrac{63 : 3}{150 : 3} = \dfrac{21}{50}.
\]
Nhân cả tử số và mẫu số với $2$:
\[
\dfrac{21}{50} = \dfrac{21 \times 2}{50 \times 2} = \dfrac{42}{100}.
\]
Vậy tử số của phân số thập phân đó là $42$.
}
\end{ex}

\begin{ex} %[5D1V4]
Tìm số tự nhiên $y$ biết: $\dfrac{y}{200} = \dfrac{35}{1\,000}$.
\shortans{7}
\loigiai{
\textbf{Đáp số:} $7$\\
Ta có:
\[
\dfrac{35}{1\,000} = \dfrac{35 : 5}{1\,000 : 5} = \dfrac{7}{200}.
\]
Vậy $y = 7$.
}
\end{ex}

\begin{ex} %[5D1V4]
Một trang trại nuôi $250$ con gà, trong đó có $85$ con gà trống. Hãy tìm tử số của phân số thập phân có mẫu số là $1\,000$ biểu thị số gà trống so với tổng số gà của trang trại.
\shortans{340}
\loigiai{
\textbf{Đáp số:} $340$\\
Phân số chỉ số gà trống là $\dfrac{85}{250}$.\\
Nhân cả tử số và mẫu số với $4$:
\[
\dfrac{85}{250} = \dfrac{85 \times 4}{250 \times 4} = \dfrac{340}{1\,000}.
\]
Tử số cần tìm là $340$.
}
\end{ex}

\begin{ex} %[5D1V4]
Viết đại lượng $48\text{ cm}$ dưới dạng phân số thập phân có đơn vị mét với mẫu số là $100$. Tử số của phân số thập phân đó là bao nhiêu?
\shortans{48}
\loigiai{
\textbf{Đáp số:} $48$\\
Vì $1\text{ m} = 100\text{ cm}$ nên $48\text{ cm} = \dfrac{48}{100}\text{ m}$.\\
Tử số của phân số đó là $48$.
}
\end{ex}

\begin{ex} %[5D1V4]
Một khu đất có diện tích $500\text{ m}^2$, người ta dùng $165\text{ m}^2$ để xây nhà. Tìm tử số của phân số thập phân có mẫu số là $100$ biểu thị diện tích xây nhà so với diện tích khu đất.
\shortans{33}
\loigiai{
\textbf{Đáp số:} $33$\\
Phân số biểu thị diện tích xây nhà là $\dfrac{165}{500}$.\\
Chia cả tử số và mẫu số cho $5$:
\[
\dfrac{165 : 5}{500 : 5} = \dfrac{33}{100}.
\]
Tử số cần tìm là $33$.
}
\end{ex}

\begin{ex} %[5D1V4]
Tìm số tự nhiên $a$ biết rằng phân số $\dfrac{a}{40}$ khi chuyển thành phân số thập phân có mẫu số là $100$ thì có tử số bằng $65$.
\shortans{26}
\loigiai{
\textbf{Đáp số:} $26$\\
Theo đề bài:
\[
\dfrac{a}{40} = \dfrac{65}{100}.
\]
Rút gọn phân số $\dfrac{65}{100}$ bằng cách chia cả tử và mẫu cho $5$:
\[
\dfrac{65}{100} = \dfrac{13}{20}.
\]
Quy đồng mẫu số về $40$:
\[
\dfrac{13}{20} = \dfrac{13 \times 2}{20 \times 2} = \dfrac{26}{40}.
\]
Do đó $a = 26$.
}
\end{ex}
